% !TEX root = ../main.tex

\phantomsection
\chapter*{초록}

\label{ch:abstract}

디파이(DeFi, 탈중앙 금융) 서비스들은 블록체인에 남은 지갑의 평판을 보고 거래 상대를 가려내고, 위험을 막는 기준을 정하고, 자동 프로그램이나 나쁜 행동을 찾아내요. 이미 발표된 그래프 방법인 AWP(적응형 가중 페이지랭크)는 지갑들이 토큰을 주고받은 송금 그래프(점과 화살표로 이은 그림) 위에서 지갑에 순위를 매겨요. 하지만 송금은 돈을 냈다는 기록일 뿐, 믿는다는 기록은 아니에요. 모든 ERC-20 토큰(토큰 공통 규칙을 따르는 토큰)에는 더 분명한 행동이 하나 더 기록돼요. 바로 허락(allowance)이에요. 허락은 토큰 주인이 스펜더(spender, 허락을 받아 대신 쓰는 쪽)에게 “내 토큰을 대신 옮겨도 돼”라고 정해 주는 일이에요. 이 논문은 허락 화살표가 평판 점수인 페이지랭크(PageRank)에 무엇을 보태는지 물어요. 그 점수는 무엇을 재는지, 점수를 만든 화살표를 그냥 센 수보다 나중 행동을 더 잘 미리 맞히는지, 계산하는 데 얼마나 드는지, 그리고 두 종류의 화살표를 한 번에 걸으면 더 나은지를 물어요.

엔도스랭크(EndorseRank)는 0이 아닌 마지막 허락들로 만든 그래프에서 계산한 페이지랭크예요. 화살표 무게는 AWP와 같은 방법으로 정하고, 모든 점에서 고르게 다시 출발해요. 엔도스랭크는 라우터나 금고(vault) 같은, 허락을 받는 스펜더에게 순위를 매겨요. 허락을 해 주는 주인에게는 순위를 매기지 않아요. 결합 페이지랭크(C-PR)는 지갑마다 허락 쪽 줄과 송금 쪽 줄을 정해진 비율 $\lambda$로 섞어서 걸어요. 씨앗을 쓰는 변형(S-PR)은 송금 그래프를 걷되, 허락 그래프를 걸은 결과에서 다시 출발해요. 수를 세어 비교하는 방법(양적 설계)으로, AWP와 원시 차수(화살표를 그냥 센 수)를 포함한 모든 점수를 똑같은 GMX~V2 트레이더 지갑 $5{,}521$개에 대해 계산했어요. 이 지갑들은 2025년 12월부터 2026년 5월까지 아비트럼 원(Arbitrum One)에서 거래했어요. 점수들은 켄달의 타우($\tau$, 두 줄 세우기가 얼마나 같은 순서인지 나타내는 수)와, 지갑을 짝지어 다시 뽑아 구한 범위(부트스트랩 신뢰구간)로 비교했어요. 가장 중요한 시험은 시간 홀드아웃(나중 기간을 떼어 두고 맞혀 보는 시험)이에요. 2026년 2월 28일에 허락을 받아 두고 있던 스펜더 $1{,}335$개를 대상으로, 그날 고정한 점수를 3월부터 5월까지 센 새 허락과 새 송금자에 견주었어요. C-PR에 대한 규칙 하나는 결합 점수를 하나라도 계산하기 전에 정해 두었어요. 두 번째 규칙은 6월부터 8월까지의 기간을 위해, 그 기간의 기록을 꺼내기 전에 미리 등록해 두었어요.

점수를 매긴 기간 밖에서는, 어떤 페이지랭크도 같은 종류 화살표의 원시 차수만큼 라벨(나중에 채점에 쓰는 정답 값)을 잘 줄 세우지 못했어요. 동결일(점수를 고정한 날)의 받은 허락 수는 나중에 허락을 새로 받은 스펜더를 $\tau=0.711$(95\% 신뢰구간 $[0.665,0.759]$)로 줄 세웠어요. 허락 그래프를 걷는 점수들이 가장 가까웠지만, 금액을 그대로 쓴 허락 층은 $0.479$, 엔도스랭크는 $0.394$였어요. 송금만으로 만든 점수들은 $0.123$ 이하에 머물렀어요. 송금 입차수(송금을 보내온 서로 다른 주소 수)는 새 송금자를 $0.472$로 줄 세웠고, AWP는 $0.344$였어요. 6월부터 8월까지의 기간에서도 두 차이가 똑같이 나타났어요. C-PR은 첫 번째 규칙을 통과하지 못했어요. 이 규칙은 C-PR이 새 허락에서는 허락 걷기만큼 잘하고($\Delta\tau=-0.184$, $[-0.213,-0.149]$), 새 송금자에서는 송금 걷기보다 잘하기($+0.007$, $[-0.012,0.028]$)를 요구했어요. 등록해 둔 기간에서 C-PR은 새 허락($+0.179$, $[0.132,0.225]$)과 트레이더의 청산 없이 닫은 비율($+0.018$, $[0.011,0.024]$)에서 송금 걷기를 이겼어요. 하지만 새 송금자에서는 범위의 아래 끝($-0.025$)이 미리 정한 허용 폭 $-0.02$보다 낮았어요. 그래서 이 논문은 두 그래프를 섞으면 도움이 된다고 주장하지 않아요. AWP와 비교하면 C-PR은 새 송금자에서 $0.073$만큼 졌어요. 같은 기간 안에서는 점수마다 자기 화살표로 센 통계와 잘 맞았어요. AWP는 송금과($0.612$), 엔도스랭크는 허락과($0.375$) 맞았어요. 이 일치는 처음부터 어느 정도 생길 수밖에 없는 것이에요. 또 엔도스랭크는 트레이더 지갑 가운데 131개만 동점인 두 값에서 떼어 낼 수 있었어요. 엔도스랭크 계산은 AWP보다 약 9배 빨랐어요($0.55$초 대 $5.03$초). 화살표 수가 9분의 1이니 그만큼 빠른 것이에요. 송금 그래프도 함께 걷는 결합 점수들은 AWP보다 더 오래 걸렸어요.

이 논문이 새로 보탠 것은 세 가지예요. 온체인(블록체인 위의) 평판을 위한 허락 화살표, 모든 점수를 그 점수가 매끄럽게 다듬는 화살표 수와 견준 기간 밖 평가, 그리고 미리 정한 규칙으로 시험하고 실패했다고 그대로 보고한 결합 걷기예요. 주인들이 어떤 스펜더를 믿는지 알고 싶은 서비스에게는 받은 허락 수가 가장 싸고 쓸모 있는 신호예요. 엔도스랭크는 그 수를 그래프 전체로 퍼뜨린 모양이에요. 둘 다 누구나 볼 수 있는 기록에서 나오므로, 누구든 다시 계산하고 확인할 수 있어요. 하지만 둘 다 허락을 해 주는 트레이더에게는 순위를 매기지 못해요. 또 모든 점에서 고르게 다시 출발하는 방식에서는 주소 10개로 만든 농장(공격자가 만든 가짜 주소 무리)이 목표 주소를 이 묶음에서 가장 높은 엔도스랭크의 7.7배까지 올릴 수 있어요. 그래서 시빌 공격(가짜 주소를 많이 만들어 속이는 공격)을 막는 장치 없이 이 점수로 대출이나 보상을 정하면 안 돼요.

\emph{핵심어---} 디파이(DeFi), 엔도스랭크(EndorseRank), 적응형 가중 페이지랭크(AWP), 결합 페이지랭크, 온체인 평판, ERC-20 허락, 시간 홀드아웃, 사전 등록, 아비트럼 원(Arbitrum One), 켄달의 타우
